%%%% Draft of 2026-10-04 (evening). NOT \input into main.tex yet. Replaces Sec. "When Recall Beats Relearning" and the
%%%% \todo in Definition 1. Chain: R-divergence for an in-context learner -> its blind spot -> excess-risk form
%%%% (= conditional Jensen-Shannon) -> merge rule = estimate of the recall regret -> meta expert under a margin
%%%% (MOOE decomposition) -> recall versus relearn. Numbers quoted here are from DEV seeds only (see the log).

\section{Theory: From Model-Oriented Discrepancy to Recall versus Relearn}
\label{sec:theory}

\textbf{Setting.} $f_\Theta(\cdot\mid\vx;\cC)$ is the predictive distribution of the frozen TFM given a context $\cC$.
For a concept $P$ over $(\vx,y)$ and a loss $\ell$, the risk of a context is
$\epsilon_P(\cC)=\E_{(\vx,y)\sim P}\,\ell\bigl(f_\Theta(\cdot\mid\vx;\cC),y\bigr)$, and the \emph{learning curve} of $P$ is
$R_P(n)=\E_{\cC\sim P^n}\,\epsilon_P(\cC)$. We write $p_\vx=P(\cdot\mid\vx)$, $P_X$ for the marginal, and
$\mathrm{JS}$ for the Jensen--Shannon divergence. Unless stated otherwise $\ell$ is the log-loss.

\subsection{A Model-Oriented Discrepancy for In-Context Learners}

R-divergence \citep{zhao2023rdiv} declares two distributions identical \emph{for a model} when the hypothesis learned
on their mixture has the same risk on both. For an in-context learner the hypothesis learned on the mixture is the
TFM conditioned on the union of the two contexts:
\begin{equation}\label{eq:rdiv}
  \Div_R(P\|Q)=\bigl|\epsilon_P(\cC_P\cup\cC_Q)-\epsilon_Q(\cC_P\cup\cC_Q)\bigr| .
\end{equation}
To study its population behaviour we use the idealisation behind prior-fitted networks
\citep{muller2022pfn,nagler2023statistical}: with contexts of equal size from $P$ and $Q$, the predictive distribution
tends to the conditional of the mixture $U=\tfrac12(P+Q)$,
$u_\vx=\alpha(\vx)p_\vx+(1-\alpha(\vx))q_\vx$ with $\alpha=P_X/(P_X+Q_X)$, and with a context from $P$ alone to
$p_\vx$. We denote these limits by $h_U$ and $h_P$.

\begin{definition}[Excess-risk discrepancy]\label{def:excess}
The \emph{recall regret} of $P$ on $Q$ and the \emph{excess-risk discrepancy} are
\begin{align}
  \Div_T(Q\,|\,P) &= \epsilon_Q(h_P)-\epsilon_Q(h_Q), \label{eq:transfer}\\
  \Div_E(P,Q) &= \tfrac12\bigl[\epsilon_P(h_U)-\epsilon_P(h_P)\bigr]+\tfrac12\bigl[\epsilon_Q(h_U)-\epsilon_Q(h_Q)\bigr].
  \label{eq:excess}
\end{align}
\end{definition}
$\Div_T$ is what recalling the stored context of $P$ costs while concept $Q$ is active; $\Div_E$ is what pooling the
two contexts costs each of them.

\begin{proposition}[Blind spot of R-divergence and its repair]\label{prop:blind}
Under the log-loss and the idealisation above,
\begin{enumerate}
  \item $\Div_E(P,Q)=\mathrm{JS}(P,Q)-\mathrm{JS}(P_X,Q_X)\ge0$, the part of the Jensen--Shannon divergence between
        the joints that covariate shift does not explain. If $P_X=Q_X$ then
        $\Div_E(P,Q)=\E_{\vx}\,\mathrm{JS}(p_\vx,q_\vx)$, which is zero if and only if $p_\vx=q_\vx$ almost surely,
        and $\Div_T(Q\,|\,P)=\E_{\vx}\KL(q_\vx\|p_\vx)$.
  \item $\Div_R(P\|Q)=\bigl|\E_{P_X}[H(p_\vx)+\KL(p_\vx\|u_\vx)]-\E_{Q_X}[H(q_\vx)+\KL(q_\vx\|u_\vx)]\bigr|$, with $H$
        the entropy.
  \item (Blind spot.) Let $P_X=Q_X$ and let both concepts be deterministic, $y=g_P(\vx)$ and $y=g_Q(\vx)$, with
        disagreement mass $\rho=\Pr_{\vx}[g_P(\vx)\neq g_Q(\vx)]$. Then $\Div_R(P\|Q)=0$ for every $\rho\in[0,1]$,
        whereas $\Div_E(P,Q)=\rho\ln2$. Under the 0-1 loss with ties of $u_\vx$ broken independently of the source
        of the query, $\epsilon_P(h_U)=\epsilon_Q(h_U)=\rho/2$, so again $\Div_R=0$, while
        $\Div_T(Q\,|\,P)=\rho=2\,\Div_E(P,Q)$.
\end{enumerate}
\end{proposition}

\begin{proof}
For the log-loss, $\epsilon_P(h)=\E_{P_X}\bigl[H(p_\vx)+\KL(p_\vx\|h_\vx)\bigr]$ for any conditional $h$, hence
$\epsilon_P(h_P)=\E_{P_X}H(p_\vx)$ and $\epsilon_P(h_U)-\epsilon_P(h_P)=\E_{P_X}\KL(p_\vx\|u_\vx)$, and likewise for $Q$;
this gives (2) and $\Div_T(Q\,|\,P)=\E_{Q_X}\KL(q_\vx\|p_\vx)$. By the chain rule,
$\KL(P\|U)=\KL(P_X\|U_X)+\E_{P_X}\KL(p_\vx\|u_\vx)$ and the same for $Q$; averaging the two identities gives
$\mathrm{JS}(P,Q)=\mathrm{JS}(P_X,Q_X)+\Div_E(P,Q)$, and $\Div_E\ge0$ as a mean of KL terms. If $P_X=Q_X$ then
$\alpha\equiv\tfrac12$, $u_\vx=\tfrac12(p_\vx+q_\vx)$ and $\Div_E=\E_\vx\mathrm{JS}(p_\vx,q_\vx)$, which vanishes iff
$p_\vx=q_\vx$ a.s. This proves (1). For (3), deterministic concepts have $H(p_\vx)=H(q_\vx)=0$, and
$u_\vx$ puts mass $\tfrac12$ on each of $g_P(\vx)$ and $g_Q(\vx)$ where they differ and mass $1$ where they agree, so
$\KL(p_\vx\|u_\vx)=\KL(q_\vx\|u_\vx)=\ln2\cdot\mathbf 1[g_P(\vx)\neq g_Q(\vx)]$; the two terms in (2) cancel and
$\Div_E=\rho\ln2$. Under the 0-1 loss the prediction at a disagreeing $\vx$ is $g_P(\vx)$ or $g_Q(\vx)$ with probability
$\tfrac12$ each, whichever distribution the query comes from, so each risk is $\rho/2$; $h_P$ errs on $Q$ exactly where
the concepts disagree, so $\Div_T(Q\,|\,P)=\rho$.
\end{proof}

\textbf{Reading.} R-divergence asks whether the mixture hypothesis treats the two sets \emph{equally}; two concepts
that contradict each other on the same inputs are treated equally badly, so the answer is yes although no context
of one helps the other. The excess-risk form asks whether pooling \emph{costs} anything, and for the log-loss it is
exactly the concept-shift part of the Jensen--Shannon divergence. This is the case that matters for recurring
concepts: on the controlled grid (shared $P_X$, labels reassigned per concept) the estimate of $\Div_R$ on pairs of
different concepts averages 0.05 while $\rho$ ranges from 0.13 to 1.00, and the estimates of $\Div_T$ and of
$2\Div_E$ track $\rho$.\todo{Dev seeds 0, 1: mean $\widehat{\Div}_T$ 0.695, mean $2\widehat{\Div}_E$ 0.657, correlation
0.99 over 72 pairs; AUROC for ``different concept'' 0.86--0.89 for $\Div_R$ and 1.00 for the other two. Replicate on
test seeds before quoting.}

\textbf{Finite contexts: pooling against interference.} With $n$ rows per set the empirical counterpart of
\eqref{eq:excess} contains a second term. If $P=Q$, pooling doubles the context, and
\begin{equation}\label{eq:pool}
  \E\bigl[\epsilon_P(\cC_P\cup\cC_Q)-\epsilon_P(\cC_P)\bigr]=R_P(2n)-R_P(n)\le0
\end{equation}
whenever the learning curve is non-increasing in risk: the excess is minus the learning-curve increment. In general
the empirical excess is interference (\cref{prop:blind}) minus pooling gain \eqref{eq:pool}, so the rule ``merge a
segment into a stored expert when the estimated excess is at most $\delta$'' merges exactly when pooling helps more
than the concepts interfere. \method{} applies the directed form: a closed segment $\cC_Q$ joins the stored expert
$\cC_P$ that minimises $\widehat{\Div}_T(Q\,|\,P)=\hat\epsilon_{Q}(\cC_P)-\hat\epsilon_{Q}(\cC_Q)$ if this is at most
$\delta$, with $\hat\epsilon_Q(\cC_Q)$ cross-fitted because an in-context learner memorises its context (the
overfitting argument of \citet{zhao2023rdiv} against evaluating on the training set). Thus the merge statistic is an
estimate of the regret that recalling this expert would incur on the segment, which is the quantity the next two
results are about.

\subsection{The Meta Expert Under a Margin}

Following \citet{zhao2025mooe}, the regret of the mixture against the per-batch best predictor splits into the regret of
the meta expert against the best expert $k^\star$ and the regret of that expert,
$\Regret_{\method}=\Regret_{\mathrm{ME}}+\Regret_{\mathrm{KE}}$. MOOE bounds $\Regret_{\mathrm{ME}}\le\sqrt{T\ln K}$ for
cumulative losses over an interval of $T$ rounds. After a recurrence the cumulative loss still carries the previous
concept, so \method{} restarts the meta expert every batch (weights from the last labelled batch only). The
worst-case bound is then vacuous, and the guarantee comes from a margin instead.

\begin{lemma}[Margin]\label{lem:margin}
Let $w_k\propto\exp(-\eta\,\ell_k)$ over $K$ experts and suppose $\ell_{k^\star}\le\ell_k-g$ for all $k\ne k^\star$. With
$\kappa=\ln\bigl(1+(K-1)e^{-\eta g}\bigr)$: (i) $w_{k^\star}\ge e^{-\kappa}$; (ii) for every row, the log-loss of the
mixture is at most that of expert $k^\star$ plus $\kappa$; (iii) the mixture predicts the same class as expert
$k^\star$ on every row where the top two probabilities of $k^\star$ differ by more than $e^{\kappa}-1$.
\end{lemma}
\begin{proof}
(i) $w_{k^\star}=\bigl(1+\sum_{k\ne k^\star}e^{-\eta(\ell_k-\ell_{k^\star})}\bigr)^{-1}\ge(1+(K-1)e^{-\eta g})^{-1}$.
(ii) $-\ln\sum_kw_kp_k(y)\le-\ln\bigl(w_{k^\star}p_{k^\star}(y)\bigr)\le-\ln p_{k^\star}(y)+\kappa$.
(iii) For classes $y_1,y_2$, $\sum_kw_k\bigl(p_k(y_1)-p_k(y_2)\bigr)\ge w_{k^\star}\bigl(p_{k^\star}(y_1)-p_{k^\star}(y_2)\bigr)
-(1-w_{k^\star})$, which is positive when the gap of $k^\star$ exceeds $(1-w_{k^\star})/w_{k^\star}\le e^{\kappa}-1$.
\end{proof}

\begin{proposition}[Regret of the meta expert on a block]\label{prop:block}
Consider a block of $J$ batches of one concept and log-losses clipped at $c$. Let $\mathcal G$ be the set of batches
$j\ge2$ of the block before which expert $k^\star$ led the last labelled batch by a margin of at least $g$. Then,
per row,
\[
  \Regret_{\mathrm{ME}}\;\le\;c\,(J-|\mathcal G|)+|\mathcal G|\,\ln\bigl(1+(K-1)e^{-\eta g}\bigr).
\]
If the margin holds from the second batch on, $\Regret_{\mathrm{ME}}\le c+(J-1)(K-1)e^{-\eta g}$: one batch of
identification, independent of the block length up to an exponentially small term, against $\sqrt{JB\ln K}$ for the
cumulative rule.
\end{proposition}
\begin{proof}
On a batch in $\mathcal G$ apply \cref{lem:margin}(ii) with the losses of the last labelled batch; on the other
batches the per-row regret is at most $c$. The special case uses $|\mathcal G|=J-1$ and $\ln(1+x)\le x$.
\end{proof}
\todo[inline]{Dev grid (12 streams): from the second batch of a recurring block the leader has weight 1.000 in all six
cells and the gap between the best and second-best expert is 0.8--8.0 nats, so with $\eta=10$, $K\le15$:
$(K-1)e^{-\eta g}\le 14e^{-8}<0.005$. Replaying the MOOE meta expert with intervals of one batch and its own step size
changes accuracy by 0.04 points; intervals of 500 rows lose 9.5 points on the hardest cell.}

\subsection{When Recall Beats Relearning}

Let concepts recur in blocks of $L$ rows ($J=L/B$ batches), labels arrive one batch late, the budget is $M$ rows, and
$a_k(n)$ is the accuracy learning curve of concept $k$, non-decreasing in $n$. On the $c$-th visit of concept $k$
($c\ge2$) its stored expert holds $n_c=\min(M,(c-1)L)$ rows. The \emph{relearning} learner resets its context at the
block boundary; in MOOE's terms it is the online expert, and the stored context is an offline expert that needs no
training.

\begin{proposition}[Recall versus relearn]\label{prop:recall}
Assume the margin of \cref{prop:block} holds from the second batch of the block for the better of the stored expert
and the window, with $e^\kappa-1$ below the top-two gap of that expert. Then on the $c$-th visit the expected
accuracy of \method{} minus that of relearning, averaged over the block, is
\begin{equation}\label{eq:area}
  \Delta_{k,c}(L)=\frac1J\sum_{j=1}^{J-1}\Bigl[a_k(n_c)-a_k\bigl(\min(jB,M)\bigr)\Bigr]_+ ,
\end{equation}
the area by which the learning curve lies below the accuracy of the stored expert. It is zero when
$a_k(B)=a_k(M)$, at most $\tfrac{M}{L}\bigl(a_k(M)-a_k(B)\bigr)$, and zero on the first visit.
\end{proposition}
\begin{proof}
Both learners lose the first batch, whose labels have not arrived. Before batch $j+1$ the relearning context holds
$\min(jB,M)$ rows of the block, so its expected accuracy is $a_k(\min(jB,M))$; the stored expert has expected accuracy
$a_k(n_c)$. By \cref{lem:margin}(iii) the mixture predicts as the better of the two, which gives the positive part.
Each term is at most $a_k(M)-a_k(B)$ and vanishes for $jB\ge M$, so at most $M/B$ of the $J$ terms are non-zero.
\end{proof}

\textbf{Reading.} \Cref{eq:area} is $\Regret_{\mathrm{OE}}-\Regret_{\mathrm{KE}}$ in the decomposition of
\citet{zhao2025mooe}: memory helps when the area between the learning curve and the stored accuracy exceeds what the
meta expert costs, and \cref{prop:block} makes that cost one batch, which the relearning learner loses as well. The
value of memory is therefore a property of the backbone and the concept that can be measured before deployment. Two
caveats: the stored expert keeps growing during the block, so \eqref{eq:area} is a lower bound once $n_c<M$; and the
relearning learner is given the block boundary, so a real detector does worse.
